% !TeX program = xelatex
\documentclass[11pt]{article}
% Shared typography and author block; build with XeLaTeX, BibTeX, XeLaTeX twice.
\usepackage[a4paper,margin=25mm]{geometry}
\usepackage{amsmath,amssymb,amsthm,mathtools,bm}
\usepackage{booktabs,array,tabularx,longtable}
\usepackage{cite}
\usepackage[hidelinks,unicode]{hyperref}
\usepackage{xurl,authblk,enumitem}
\setlist[enumerate]{leftmargin=2em,itemsep=0.3ex,topsep=0.5ex}
\setlength{\emergencystretch}{3em}
\linespread{1.08}
\clubpenalty=10000
\widowpenalty=10000
\displaywidowpenalty=10000
\numberwithin{equation}{section}
\newtheorem{theorem}{Theorem}[section]
\newtheorem{lemma}[theorem]{Lemma}
\newtheorem{proposition}[theorem]{Proposition}
\newtheorem{corollary}[theorem]{Corollary}
\theoremstyle{definition}
\newtheorem{definition}[theorem]{Definition}
\theoremstyle{remark}
\newtheorem{remark}[theorem]{Remark}
\newcommand{\R}{\mathbb R}
\newcommand{\C}{\mathbb C}
\newcommand{\Repart}{\operatorname{Re}}
\newcommand{\conv}{\operatorname{conv}}
\newcommand{\SPR}{\mathrm{SPR}}
\newcommand{\code}[1]{\texttt{\small\detokenize{#1}}}
\newcommand{\id}[1]{\nolinkurl{#1}}
\newcommand{\doi}[1]{doi: \href{https://doi.org/#1}{\nolinkurl{#1}}}
\renewcommand{\Authfont}{\small}
\renewcommand{\Affilfont}{\small}
\renewcommand{\Authsep}{, }
\renewcommand{\Authand}{, }
\renewcommand{\Authands}{, }
\author[1]{Zhengtai Xiong}
\author[1]{Tiancheng Xu}
\author[1]{Xiang Ji}
\author[1]{Guowei Dou}
\author[2,3]{Long Wang}
\author[1]{Wensheng Yu}
\affil[1]{School of Electronic Engineering, Beijing University of Posts and Telecommunications, Beijing 100876, China}
\affil[2]{Center for Systems and Control, College of Engineering, Peking University, Beijing 100871, China}
\affil[3]{Center for Multi-Agent Research, Institute for Artificial Intelligence, Peking University, Beijing 100871, China}
\date{}

\title{Formalizing the Yu--Wang Quartic SPR Synthesis Theorem with a Single Perturbation Parameter}
\hypersetup{pdftitle={Formalizing the Yu--Wang Quartic SPR Synthesis Theorem with a Single Perturbation Parameter}}
\begin{document}
\maketitle
\begin{abstract}
We formalize the Yu--Wang SPR synthesis theorem for fourth-order polynomial segments in Lean. The established criterion states that two monic quartic denominators admit a common equal-degree strictly positive real numerator if and only if their whole convex segment is Hurwitz stable. The proof retains the original elliptic-region strategy. Explicit support functions establish intersection of the two regions, and a completed-square margin permits a single positive parameter to perturb the constant and leading coefficients simultaneously. The numerator can be chosen as $\varepsilon s^4+s^3+xs^2+ys+\varepsilon$, with an explicit admissible bound on $\varepsilon$. The converse is proved independently by an interlacing argument for numerator stability and by reciprocation. The Lean development imports only Mathlib and proves the theorem for the explicit quartic Routh--Hurwitz coefficient predicate. The contribution is a streamlined formalization with a quantitative perturbation step, not a new segment criterion.
\end{abstract}
\noindent\textbf{Keywords:} quartic polynomial segment; strictly positive real synthesis; formal verification; ellipse; perturbation.

\section{Introduction and attribution}
Following the cubic synthesis result of Yu and Huang \cite{yuHuang1999}, Yu and Wang proved that robust Hurwitz stability of a quartic polynomial segment is equivalent to the existence of a common equal-degree SPR numerator \cite{yuWang2001Quartic}. Their proof constructs an open ellipse for each denominator, proves intersection through a separating tangent line, obtains a cubic numerator by a small constant perturbation, and then raises its degree. This geometric route is the mathematical foundation of the present formalization.

The file \code{SPR/simplified/n4.lean} retains the ellipse and separation mechanism. It replaces tangent-line elimination by a dual quadratic identity, and it chooses the constant perturbation and leading coefficient to be the same positive number. These changes provide explicit inequalities convenient for a proof assistant. They do not change the original synthesis criterion or replace its principal geometric idea. Accordingly, this article is positioned as a formalization and refinement of \cite{yuWang2001Quartic}.

The supplied arbitrary-order manuscript \cite{referenceDraft} guides the presentation, but its finite-contact proof is not invoked. The quartic Lean file is independent of the cubic, quintic, sextic, and arbitrary-order synthesis files. In particular, the existence theorem is proved within the quartic coefficient model rather than imported as a specialization of a more general result.

\section{The coefficient model and synthesis statement}
Let
\[
 a(s)=s^4+a_1s^3+a_2s^2+a_3s+a_4,\qquad
 b(s)=s^4+b_1s^3+b_2s^2+b_3s+b_4,
 \qquad p_\lambda=(1-\lambda)a+\lambda b.
\]
Define $H_4(a)$ by
\begin{equation}\label{eq:rh}
 a_j>0\ (1\le j\le4),\qquad a_1a_2-a_3>0,\qquad
 \Delta_a:=a_1a_2a_3-a_3^2-a_1^2a_4>0.
\end{equation}
This is the classical quartic Routh--Hurwitz coefficient condition. In the formal file it is the definition of \code{Hurwitz}. The usual interpretation in terms of roots in the open left half-plane is external to this coefficient-model formalization; the file does not prove a general Routh--Hurwitz theorem. Notice that positivity of all coefficients and $\Delta_a>0$ already imply $a_1a_2-a_3>0$.

For a real quartic $c$ with positive leading coefficient, call $c/a$ SPR when $H_4(a)$ holds and
\begin{equation}\label{eq:spr}
 \Repart\frac{c(i\omega)}{a(i\omega)}>0\quad(\omega\in\R).
\end{equation}
Put
\[
 E_a(t)=t^2-a_2t+a_4,\qquad O_a(t)=a_3-a_1t.
\]
Then $a(i\omega)=E_a(\omega^2)+i\omega O_a(\omega^2)$. Since
\begin{equation}\label{eq:pivot}
 a_1^2E_a(a_3/a_1)=-\Delta_a<0,
\end{equation}
the even and odd parts cannot vanish together on the imaginary axis; $a_4>0$ handles $\omega=0$. Thus the quotient in \eqref{eq:spr} is defined at every real frequency. Equal-degree monic quotients tend to $1$ at infinity.

\begin{theorem}[Formalized quartic segment criterion]\label{thm:main}
Assume $H_4(a)$ and $H_4(b)$. The following are equivalent:
\begin{enumerate}
\item $H_4(p_\lambda)$ for every $\lambda\in[0,1]$;
\item a common monic quartic SPR numerator exists;
\item a common quartic SPR numerator with positive leading coefficient exists.
\end{enumerate}
If these conditions hold, there are $x,y\in\R$ and $\varepsilon>0$ such that
\begin{equation}\label{eq:c}
 c_\varepsilon(s)=\varepsilon s^4+s^3+xs^2+ys+\varepsilon
\end{equation}
works at both endpoints. The cubic $q_\varepsilon=s^3+xs^2+ys+\varepsilon$ also has strictly positive real part after division by either endpoint at every real frequency.
\end{theorem}
The reduced-degree statement is stated separately because the equal-degree SPR convention used here does not call $q_\varepsilon/a$ a quartic SPR quotient.

\section{The ellipse and its support function}
For $x,y\in\R$ define
\begin{align}
 F_a(t)&=(a_1-x)t^2+(a_2x-a_1y-a_3)t+(a_3y-a_4x),\label{eq:F}\\
 C_a(x,y)&=(a_2x-a_1y-a_3)^2-4(a_1-x)(a_3y-a_4x),\label{eq:conic}\\
 \Omega_a&=\{(x,y):C_a(x,y)<0\}.
\end{align}
Expansion of \eqref{eq:conic} yields
\begin{align*}
 C_a(x,y)={}&(a_2^2-4a_4)x^2+2(2a_3-a_1a_2)xy+a_1^2y^2\\
 &-2(a_2a_3-2a_1a_4)x-2a_1a_3y+a_3^2,
\end{align*}
the ellipse polynomial of \cite[Lemma~1]{yuWang2001Quartic}. Its center is $(a_1/2,a_2/2)$. With $X=x-a_1/2$, $Y=y-a_2/2$,
\[
 C_a(x,y)=H_a(X,Y)-\Delta_a,
\]
where
\[
 H_a(X,Y)=(a_2^2-4a_4)X^2+2(2a_3-a_1a_2)XY+a_1^2Y^2.
\]
The identity
\begin{equation}\label{eq:square}
 a_1^2H_a(X,Y)
 =\bigl(a_1^2Y+(2a_3-a_1a_2)X\bigr)^2+4\Delta_aX^2
\end{equation}
shows that $H_a$ is positive definite and $\Omega_a$ is a nonempty open convex ellipse. It also gives $|X|<a_1/2$ on $\Omega_a$, so $a_1-x>0$. Completing the square in \eqref{eq:F} gives
\begin{equation}\label{eq:margin}
 F_a(t)\ge m_a:=\frac{-C_a(x,y)}{4(a_1-x)}>0\qquad(t\in\R).
\end{equation}

For a nonzero normal $(A,B)$ introduce
\[
 D_a(A,B)=a_1^2A^2-2(2a_3-a_1a_2)AB+(a_2^2-4a_4)B^2,
 \qquad r_a=\sqrt{D_a(A,B)}>0.
\]
The determinant of the matrix of $H_a$ is $4\Delta_a$. Consequently the support values of the open ellipse are
\begin{equation}\label{eq:support}
 \sup_{\Omega_a}(Ax+By)=\frac{a_1A+a_2B+r_a}{2},\qquad
 \inf_{\Omega_a}(Ax+By)=\frac{a_1A+a_2B-r_a}{2}.
\end{equation}
The Lean proof constructs boundary support points and approaches them from the center, so it does not assume that a supremum is attained inside an open ellipse.

\begin{lemma}\label{lem:intersection}
If the entire quartic segment satisfies $H_4$, then $\Omega_a\cap\Omega_b\ne\varnothing$.
\end{lemma}
\begin{proof}
Otherwise separation and \eqref{eq:support} provide a nonzero normal with
\[
 r_a+r_b\le L:=(b_1-a_1)A+(b_2-a_2)B.
\]
Direct expansion, including the term $-4a_4B^2$, gives
\begin{equation}\label{eq:dualblend}
 D_{p_\lambda}(A,B)
 =(1-\lambda)D_a(A,B)+\lambda D_b(A,B)-\lambda(1-\lambda)L^2.
\end{equation}
The fourth coefficient contributes affinely, so no additional quadratic remainder occurs. At $\lambda_*=r_a/(r_a+r_b)$, the weighted sum of the first two terms equals $r_ar_b$, while $\lambda_*(1-\lambda_*)L^2\ge r_ar_b$. Thus $D_{p_{\lambda_*}}(A,B)\le0$, contradicting the positive definiteness supplied by $H_4(p_{\lambda_*})$.
\end{proof}

This replaces the common-tangent coordinate calculation of \cite[Lemma~2]{yuWang2001Quartic}. Both arguments pass from disjoint elliptic regions to failure of stability at a convexly interpolated denominator; the proof strategy remains the same.

\section{A single perturbation parameter}
Choose $(x,y)\in\Omega_a\cap\Omega_b$. For a general numerator
\[
 c(s)=r s^4+s^3+xs^2+ys+z
\]
the cross polynomial is
\begin{equation}\label{eq:cross}
 Q_{a,c}(t)=tF_a(t)+zE_a(t)+rt^2E_a(t),\qquad
 \Repart\frac{c(i\omega)}{a(i\omega)}=\frac{Q_{a,c}(\omega^2)}{|a(i\omega)|^2}.
\end{equation}
The following estimate removes the need for two successive smallness choices.

\begin{lemma}\label{lem:perturb}
For $t\ge0$,
\begin{equation}\label{eq:bounds}
 E_a(t)\ge a_4-a_2t,\qquad t^2E_a(t)\ge-a_2^3t.
\end{equation}
If $r,z\ge0$ and $za_2+ra_2^3\le m_a$, then $Q_{a,c}(t)\ge za_4$ for all $t\ge0$.
\end{lemma}
\begin{proof}
The first inequality follows from $t^2\ge0$. For the second, if $t\ge a_2$ then $E_a(t)=t(t-a_2)+a_4>0$. If $0\le t\le a_2$, then
\[
 tE_a(t)=t^3-a_2t^2+a_4t\ge-a_2t^2\ge-a_2^3.
\]
Multiplication by $t\ge0$ proves the second bound. Combining these estimates with \eqref{eq:margin} in \eqref{eq:cross} yields
\[
 Q_{a,c}(t)\ge za_4+(m_a-za_2-ra_2^3)t\ge za_4.
\]
\end{proof}

Set
\begin{equation}\label{eq:epsilon}
 \varepsilon=
 \min\left\{\frac{m_a}{a_2+a_2^3},\frac{m_b}{b_2+b_2^3}\right\}>0.
\end{equation}
Both denominators in \eqref{eq:epsilon} are positive. Taking $r=z=\varepsilon$ gives
\[
 Q_{a,c_\varepsilon}(t)\ge\varepsilon a_4>0,\qquad
 Q_{b,c_\varepsilon}(t)\ge\varepsilon b_4>0\quad(t\ge0).
\]
Thus \eqref{eq:c} is a common quartic SPR numerator. Equality in the smallness bound is allowed: the strictly positive constant term $\varepsilon a_4$ remains. The same $\varepsilon$, with $r=0$, gives the common cubic $q_\varepsilon$. Finally, $c_\varepsilon/\varepsilon$ is the required monic quartic.

\begin{remark}
Formula \eqref{eq:epsilon} is an admissible bound after choosing a common ellipse point. It is not claimed to maximize a passivity margin or to be invariant under a change of frequency units. The proof establishes feasibility for the fixed polynomial coordinates of the theorem.
\end{remark}

\section{Necessity by a numerator-stability argument}
\begin{lemma}\label{lem:numstable}
If $H_4(a)$ and a monic quartic $n$ has $\Repart(n(i\omega)/a(i\omega))>0$ at every real frequency, then $H_4(n)$.
\end{lemma}
\begin{proof}
By \eqref{eq:pivot}, the monic quadratic $E_a$ has two positive roots $\alpha,\beta$ satisfying
\[
 0<\alpha<a_3/a_1<\beta.
\]
At these roots the cross polynomial reduces to $tO_n(t)O_a(t)$. Strict positivity implies
\[
 n_3-n_1\alpha>0,\qquad n_3-n_1\beta<0.
\]
Therefore $n_1>0$, $n_3>0$, and $u=n_3/n_1\in(\alpha,\beta)$. At zero frequency, positivity gives $n_4>0$. Since $E_a(u)<0$ and $O_n(u)=0$, cross positivity at $u$ forces
\[
 E_n(u)=u^2-n_2u+n_4<0.
\]
This implies $n_2>0$ and, after multiplication by $n_1^2$,
\[
 n_1n_2n_3-n_3^2-n_1^2n_4>0.
\]
Together with coefficient positivity these are exactly \eqref{eq:rh} for $n$.
\end{proof}

If a common monic numerator $n$ exists, Lemma~\ref{lem:numstable} makes $n$ a stable denominator. Reciprocation preserves strict positivity of the real part, since
\[
 \Repart\frac1z=\frac{\Repart z}{|z|^2}\quad(z\ne0).
\]
Consequently $a/n$ and $b/n$ have positive real parts, and so does $p_\lambda/n$ by linearity. Applying Lemma~\ref{lem:numstable} again proves $H_4(p_\lambda)$ for every $\lambda$. Arbitrary positive leading coefficients reduce to the monic case by scaling. This completes Theorem~\ref{thm:main}.

\begin{corollary}
The same numerator works for every denominator $p_\lambda$ on the segment.
\end{corollary}
\begin{proof}
The cross product $\Repart(c(i\omega)\overline{p_\lambda(i\omega)})$ is the convex combination of the two strictly positive endpoint cross products. Divide by $|p_\lambda(i\omega)|^2>0$.
\end{proof}

\section{Formal verification}
The proof presented in this paper has been formally verified in Lean. The code repository is \url{https://github.com/XPiPiT/n_random/tree/main/SPR/simplified}.

\section{Conclusion}
The quartic synthesis theorem admits a compact formalization that preserves the original ellipse geometry. A dual support identity proves intersection, and the inequality $t^2E_a(t)\ge-a_2^3t$ gives a simultaneous perturbation of the constant and leading coefficients. The result supplies both reduced-degree and equal-degree witnesses, together with an independent formal necessity proof.

\bibliographystyle{IEEEtran}
\bibliography{spr_references}
\end{document}
